
\documentclass[../../main.tex]{subfiles}
\begin{document}

% 年度の大きな表示
\begin{center}
  {\fontsize{48pt}{58pt}\selectfont \textbf{1962}\normalsize\textbf{年度}}
\end{center}

\vspace{10mm}

% 本文

\vspace{15mm}

% 出題テーマと難易度の表
\noindent\textbf{出題テーマと難易度}

\vspace{5mm}

\begin{center}
  \shadowbox{%
    \begin{tabular}{|c|p{8cm}|c|}
      \hline
      \textbf{問題番号} & \textbf{テーマ} & \textbf{難易度} \\
      \hline
      第1問 & 整数問題 & \\
      \hline
      第2問 & 図形と三角比 & \\
      \hline
      第3問 & 円と軌跡 & \\
      \hline
      第4問 & 三角関数の最大最小 & \\
      \hline
      第5問 & 数列と積分の極限 & \\
      \hline
      第6問 & 微分と証明 & \\
      \hline
    \end{tabular}%
  }
\end{center}

\vspace{15mm}

% 解答
\noindent\textbf{解答}

\vspace{5mm}

\begin{center}
  \shadowbox{%
    \renewcommand{\arraystretch}{1.6}
    \begin{tabular}{|c|c|p{8cm}|}
      \hline
      \textbf{問題番号} & \textbf{小問} & \textbf{答え} \\
      \hline
      第1問 & - & $1,2,3,4,6,7,8,9,12,13,14,18,19,24$ \\
      \hline
      第2問 & - & $l \fallingdotseq 6.5$（cm） \\
      \hline
      \multirow{2}{*}{第3問} & 前半 & 証明問題（定点$(3,4),(5,0)$） \\
      \cline{2-3}
                             & 後半 & $a=\dfrac{2}{3},\ 2$ \\
      \hline
      第4問 & - & $\max f=2\sqrt{2}\ \left((x,y)=\left(\dfrac{\pi}{4},\dfrac{\pi}{3}\right),\left(\dfrac{3\pi}{4},\dfrac{4\pi}{3}\right)\right)$ \\
      \hline
      第5問 & - & $\dfrac{h}{h-\sin h}\cos x$ \\
      \hline
      第6問 & - & 証明問題 \\
      \hline
    \end{tabular}%
    \renewcommand{\arraystretch}{1.0}
  }
\end{center}

\vspace{15mm}

% 解答時間
\noindent\textbf{試験形態}

\vspace{5mm}

\begin{center}
  \shadowbox{%
    \begin{tabular}{p{8cm}}
      （要確認）
    \end{tabular}%
  }
\end{center}

\end{document}
